\section*{Exercises}
\addcontentsline{toc}{section}{Exercises}

\subsection*{A. Theory}

\begin{exercise}
\exerciseid{14.A1}
Explain why $A^TA$ is symmetric and PSD for any real matrix $A$.
\end{exercise}

\begin{exercise}
\exerciseid{14.A2}
Explain, without proving the full Eckart--Young theorem, why a low-rank
approximation should keep the largest singular values first.
\end{exercise}

\subsection*{B. By Hand}

\begin{exercise}
\exerciseid{14.B1}
Find an SVD of the diagonal matrix
\[
  \begin{bmatrix}3&0\\0&1\end{bmatrix}.
\]
\end{exercise}

\begin{exercise}
\exerciseid{14.B2}
For singular values $8,3,1$, compute the squared Frobenius-norm error of the
best rank-$1$ approximation and of the best rank-$2$ approximation.
\end{exercise}

\subsection*{C. Python / NumPy}

\begin{exercise}
\exerciseid{14.C1}
Use \texttt{np.linalg.svd} to reconstruct a random $4\times 3$ matrix.
\end{exercise}

\begin{exercise}
\exerciseid{14.C2}
Center a small data matrix by subtracting column means, compute its SVD, and
report how much variance is captured by the first singular direction.
\end{exercise}

\begin{exercise}
\exerciseid{14.C3}
For a random $6\times 4$ matrix, compute rank-$k$ approximations for
$k=1,2,3$. Compare the measured Frobenius error with the formula using
discarded singular values.
\end{exercise}

\subsection*{D. Challenge}

\begin{exercise}
\exerciseid{14.D1}
Explain what singular values measure geometrically for a $2\times 2$ matrix.
\end{exercise}
